We organise the study into three tiers so that each claim is tested where it can be
tested most sharply. The \emph{simulation tier} draws pair signs directly from their
generating distribution---exactly the object Theorems~\ref{thm:null} and
\ref{thm:cert} control---which lets us calibrate validity and power over thousands of
seeds without confounding them with the vagaries of a particular model. The
\emph{end-to-end tier} fine-tunes real models on corpora with injected canaries,
unlearns them with each subject method, and certifies the result, so that the abstract
null meets an actual training pipeline. The \emph{benchmark tier} repeats the
end-to-end protocol on the field's standard testbeds, TOFU and MUSE-News, across three
architectures.

Unless noted we use $\alpha = 0.05$, two-sided certification, the sign-based
revocation arm, the finite-cohort certificate target, LoRA fine-tuning of rank $16$,
$Q = 2$ query wrappers, and between $384$ and $640$ canary pairs; the full
hyperparameter table is Appendix~\ref{app:repro}, and the headline tolerance is
$\eps = 0.2$ with the complete sweep over $\eps \in \{0.05, 0.1, 0.2\}$ reported in
Table~\ref{tab:epssweep}. The entire study ran on free, ephemeral compute---Colab and
Kaggle GPU sessions, Colab CPU sessions, an Apple-silicon laptop, and a four-core
container---coordinated by a released fault-tolerant scheduler; because every training
stage checkpoints and every verification is saved, a reclaimed session costs at most a
hundred optimiser steps.

\subsection{Which null is actually controlled}\label{sec:exp-dependence}

We begin with the question that reorganised this version of the paper. Signs from all
canaries pass through one jointly trained and unlearned model, so they are neither
independent nor identically distributed across templates and repetition strata. Does
that matter? Table~\ref{tab:dependence} answers it directly. We generate sign streams
whose \emph{marginal} sign probability is exactly $p_0 = 0.55$ under four regimes and
ask how often each certificate process issues.

The answer is that it matters, and in the direction least convenient for the earlier
analysis. When signs are independent, the fixed-boundary process of the
earlier draft holds its level at $0.034$, as its stated theorem claims. Heterogeneity
across strata alone does not break it either ($0.026$), because shuffling the reveal
order leaves the realised cohort mean tightly concentrated. But once the model itself
carries a shared random effect---so that a given training run either leaks a lot or
leaks little---the realised cohort advantage spreads out (its standard deviation rises
from $0.031$ to $0.50$) and the fixed-boundary process issues under the marginal null
in $38.9\%$ and $48.1\%$ of runs, eight- and ten-fold breaches of $\alpha = 0.05$. The
without-replacement process, by contrast, controls the realised-cohort error at or
below $0.020$ in every regime, exactly as Theorem~\ref{thm:cert} requires.

Two readings of that table are available and we want to be clear which one we take.
One could say the marginal null is the target and the procedure fails; we say instead
that the marginal null is the wrong target, because it averages over training runs the
auditor is not auditing, and a certificate should concern the deployed model. The
column marked $\dagger$ makes the distinction concrete: those issuances are not errors
under the cohort claim, because on those runs the realised advantage genuinely sits
below the tolerance. What we do concede without qualification is that the earlier
draft claimed the marginal statement while proving the conditional one, and that the
gap is not academic.

Table~\ref{tab:cohortnull} verifies the corrected theorem directly, on fixed sign
vectors sitting at the least favourable feasible point of the cohort null and carrying
deliberately adversarial internal structure---all successes contiguous, spread evenly,
or split into four strata with rates from $0.2$ to $0.9$. Under peeking after every
pair the realised error never exceeds $0.035$ against a nominal $0.05$, for every
structure, both tolerances and both cohort sizes. This is the guarantee that needs no
assumption at all.

Is the dependence real on real models? Yes. Across the benchmark runs a chi-square test
of sign-rate homogeneity across repetition strata rejects at the $5\%$ level in $16\%$
of all runs---and in $23$ of the $27$ un-unlearned runs, where there is leakage to be
heterogeneous about. Heterogeneity across canary templates rejects in $8\%$. The
regime of Table~\ref{tab:dependence} is therefore the empirically relevant one.

\begin{table}[t]
\centering\small
\caption{Which null is controlled. Sign streams with \emph{marginal} mean
$p_0 = 0.55$ ($\eps = 0.1$), $1{,}024$ pairs, $2{,}000$ seeds, peeking after every
pair. ``marginal'' is the issuance rate under the marginal null; ``cohort'' restricts
to runs whose realised cohort advantage genuinely violates the tolerance, the event
Theorem~\ref{thm:cert} controls. Bold marks a breach of $\alpha = 0.05$.
$^{\dagger}$These issuances are not errors under the cohort claim.}
\label{tab:dependence}
\begin{tabular}{lccccc}
\toprule
& & \multicolumn{2}{c}{fixed boundary (v1)} & \multicolumn{2}{c}{without replacement (ours)}\\
\cmidrule(lr){3-4}\cmidrule(lr){5-6}
sign-generating regime & $\mathrm{sd}(\hat\Delta)$ & marginal & cohort & marginal & cohort\\
\midrule
independent & 0.031 & 0.034 & 0.012 & 0.501$^{\dagger}$ & 0.014\\
heterogeneous by stratum & 0.030 & 0.026 & 0.016 & 0.495$^{\dagger}$ & 0.020\\
model-level over-dispersion & 0.296 & \textbf{0.389} & 0.001 & 0.483$^{\dagger}$ & 0.001\\
shared model-level latent & 0.500 & \textbf{0.481} & 0.000 & 0.481$^{\dagger}$ & 0.000\\
\bottomrule
\end{tabular}

\end{table}

\begin{table}[t]
\centering\small
\caption{Direct verification of Theorem~\ref{thm:cert}: realised type-I error on fixed
sign vectors at the least favourable feasible point of the cohort null, revealed in a
uniformly random order with peeking after every pair ($2{,}000$ seeds,
$\alpha = 0.05$). ``ours'' is the without-replacement process, ``fixed'' the
fixed-boundary one.}
\label{tab:cohortnull}
\begin{tabular}{llcccc}
\toprule
& & \multicolumn{2}{c}{$m=384$} & \multicolumn{2}{c}{$m=1024$}\\
\cmidrule(lr){3-4}\cmidrule(lr){5-6}
$\varepsilon$ & cohort structure & ours & fixed & ours & fixed\\
\midrule
0.1 & contiguous & 0.028 & 0.018 & 0.024 & 0.014\\
 & spread & 0.025 & 0.015 & 0.033 & 0.019\\
 & stratified (rates $0.2$--$0.9$) & 0.024 & 0.015 & 0.035 & 0.018\\
0.2 & contiguous & 0.026 & 0.013 & 0.026 & 0.013\\
 & spread & 0.026 & 0.010 & 0.034 & 0.018\\
 & stratified (rates $0.2$--$0.9$) & 0.025 & 0.013 & 0.032 & 0.015\\
\bottomrule
\end{tabular}

\end{table}

\subsection{Validity under adversarial optional stopping}\label{sec:exp-validity}

With the target settled we can ask whether VOUCH holds its error level when an
adversary peeks after every observation and stops at the first crossing. We simulate
exactly this attack. Table~\ref{tab:validity} reports the realised type-I error over
$2{,}000$ seeds: VOUCH stays below the nominal level for both the certificate and the
revocation arm at both $\alpha=0.05$ and $\alpha=0.01$, whereas a fixed-sample binomial
test used in the same peeking fashion inflates its error to $0.35$, seven times the
nominal rate. The calibration curve in Figure~\ref{fig:validity} confirms that the
conservatism holds across the whole range of $\alpha$. When we run the complete
protocol---three correlated scores, the confidence sequence, and both arms
together---the false revocation rate is $0.025$ and the per-score confidence sequence
covers the truth with uniform-in-time probability $0.971$ against the nominal $0.95$;
and when there is genuine memorisation to find ($\Delta \approx 0.36$), the revocation
arm fires in every run at a median of $44$ pairs, with no false certificate ever
issued.

One caveat on reading these numbers, which a reviewer rightly pressed. A realised error
of $0.031$ against a nominal $0.05$ is what a supermartingale bound is \emph{supposed}
to produce; it is evidence of validity, not of tightness, and the same conservatism
costs power, which is the price Table~\ref{tab:baselines} quantifies. We do not claim
the bound is tight and the gap is not a defect.

\begin{table}[t]
\centering\small
\caption{Type-I error over $2{,}000$ seeds with peeking after every pair ($1{,}024$
pairs). Bold marks a procedure whose realised error exceeds its nominal level.}
\label{tab:validity}
\begin{tabular}{lcc}
\toprule
procedure & $\alpha=0.05$ & $\alpha=0.01$\\
\midrule
VOUCH certificate (boundary null $\Delta=\varepsilon$) & 0.031 & 0.008\\
VOUCH revocation (exact unlearning) & 0.030 & 0.006\\
binomial w/ peeking --- false certification & \textbf{0.354} & \textbf{0.121}\\
binomial w/ peeking --- false alarm & \textbf{0.368} & \textbf{0.117}\\
\bottomrule
\end{tabular}

\end{table}

\begin{figure}[t]
\centering
\includegraphics[width=\linewidth]{fig1_validity.pdf}
\caption{\textbf{Anytime validity.} Left: the realised false-certification rate lies
below the nominal level at every $\alpha$ under adversarial optional stopping. Right:
a fixed-$n$ binomial test used with peeking inflates its error sevenfold; VOUCH does
not.}
\label{fig:validity}
\end{figure}

\subsection{Sound composition across scores}\label{sec:exp-soundness}

We next confirm the composition result of Theorem~\ref{thm:cert} and the failure it
warns against. Constructing a composite null in which one of three scores sits exactly
at the boundary while the other two are clean, we compare the per-score-minimum
certificate that VOUCH uses against the adaptive-mixture certificate that seems
natural. The difference is reported in Table~\ref{tab:soundness}: the mixture arm
certifies the leaking model in $99.6\%$ of runs---it is not merely loose but
anytime-invalid, as Remark~\ref{rem:mixture} explains---while the per-score minimum
holds at $0.7\%$, comfortably below $\alpha$. We want to be precise about what this
number is and is not: the mixture arm is \emph{our own} earlier design, evaluated
against an adversarially constructed composite null. It is not a published method and
the figure is not a refutation of anyone else's work. It is a trap we fell into and
climbed out of, reported so that others need not.

\begin{table}[t]
\centering\small
\caption{False certification under a composite null (score $A$ at the boundary
$p_0=0.55$, scores $B$ and $C$ clean; $2{,}048$ pairs, $\alpha=0.05$). The mixture arm
is our own earlier design, not a published method.}
\label{tab:soundness}
\begin{tabular}{lc}
\toprule
certificate arm & false-certification rate\\
\midrule
per-score minimum (ours) & 0.007\\
adaptive-mixture sign (design v1.0) & \textbf{0.996}\\
\bottomrule
\end{tabular}

\end{table}

\subsection{Comparison against valid sequential procedures}\label{sec:exp-baselines}

A fixed-sample binomial test inspected after every pair is a useful illustration of
why naive peeking is invalid, but it is not a serious competitor. The right question is
what anytime validity costs relative to procedures that are \emph{also} valid, so we
add four: exact group-sequential $\alpha$-spending with Pocock and O'Brien--Fleming
boundaries, calibrated by dynamic programming over the exact binomial lattice rather
than by a normal approximation; a truncated Beta$(\tfrac12,\tfrac12)$-mixture
(Jeffreys) e-process; a Robbins normal-mixture e-process; and a fixed-$n$ binomial test
at a pre-committed sample size. All see the same sign stream and the same query
budget. Table~\ref{tab:baselines} reports the result.

Every procedure holds its level, as it should. On power and cost the ordering is
informative. Against exact unlearning at $\eps = 0.1$ and a budget of $1{,}536$ pairs,
VOUCH on the cohort target certifies every run using a mean of $523$ queries; Pocock
reaches $0.971$ using $645$; O'Brien--Fleming $0.985$ using $844$; the two mixture
e-processes trail at $0.77$--$0.78$; and the fixed-$n$ test attains $0.989$ but spends
the entire budget of $1{,}536$ queries on every run, because a fixed-$n$ test cannot
stop early by construction. The honest summary is that anytime validity is close to
free here: against the strongest valid sequential comparator VOUCH gives up nothing in
power and saves about a fifth of the queries, and against the fixed-$n$ test it gives
up one point of power for a threefold reduction in verifier cost. The comparison also
isolates how much of VOUCH's advantage comes from retargeting: on the
super-population target, at the same budget, power falls from $1.000$ to $0.882$ and
mean queries rise from $523$ to $747$. That gap is the price of the stronger claim, and
it is the same gap Proposition~\ref{prop:transfer} charges in tolerance.

\begin{table}[t]
\centering\small
\caption{Valid sequential comparators on the same sign stream and query budget
($2{,}000$ seeds, $\alpha = 0.05$). ``type I'' is measured at the least favourable
null; for the cohort target it is restricted to runs whose realised cohort advantage
violates the tolerance. ``power'' and ``queries'' are measured under exact unlearning.}
\label{tab:baselines}
\begin{tabular}{lccccc}
\toprule
& \multicolumn{2}{c}{$\varepsilon=0.2$, $n=512$} & \multicolumn{3}{c}{$\varepsilon=0.1$, $n=1536$}\\
\cmidrule(lr){2-3}\cmidrule(lr){4-6}
procedure & type I & queries & type I & power & queries\\
\midrule
VOUCH, cohort target & 0.009 & 143 & 0.018 & 1.000 & 523\\
VOUCH, super-population target & 0.026 & 197 & 0.030 & 0.882 & 747\\
Beta$(\frac12,\frac12)$-mixture e-process & 0.033 & 224 & 0.040 & 0.769 & 900\\
normal-mixture e-process & 0.027 & 227 & 0.034 & 0.782 & 901\\
group sequential, Pocock & 0.045 & 175 & 0.050 & 0.971 & 645\\
group sequential, O'Brien--Fleming & 0.048 & 251 & 0.051 & 0.985 & 844\\
fixed-$n$ binomial at pre-committed $n$ & 0.047 & 512 & 0.053 & 0.989 & 1536\\
\bottomrule
\end{tabular}

\end{table}

\subsection{Power, tightness, and the cost of a horizon}\label{sec:exp-power}

Table~\ref{tab:power} reports certification time against the Kullback--Leibler limit.
Two corrections to the earlier version are folded in. First, the theoretical column now
follows the corrected expansion \eqref{eq:klexp}; the numbers themselves are unchanged
because they were always computed from the exact divergence, but the accompanying
sample-size rule of thumb is $2\log(1/\alpha)/\eps^2$, not $8\log(1/\alpha)/\eps^2$.
Second, and more consequentially, the earlier table's starred entries were medians
taken over the runs that happened to issue within a $20{,}000$-pair horizon, which is
a selection-biased statistic: at $\eps = 0.02$ fully $45\%$ of runs are censored, and
conditioning on the fast half produced a median of $11{,}189$ that sat \emph{below} the
information-theoretic limit of $14{,}976$, an impossibility that should have been a
red flag. Reporting the censoring rate and the Kaplan--Meier median removes the
artefact: the corrected median is $18{,}377$, comfortably above the limit, and at
$\alpha = 0.01$ censoring reaches $72\%$ and the survival curve never reaches one half,
so no median exists and we report none.

The finite-cohort column is faster than the super-population one at every tolerance,
and the reason is structural rather than statistical: the cohort claim resolves
deterministically once the cohort is exhausted, so its sample complexity is capped by
$m$ and it may legitimately sit below a bound that governs the super-population
parameter. Readers should not read the cohort column as beating an information limit;
it is answering an easier question, and Proposition~\ref{prop:transfer} prices the
difference.

Two-sided certification, which we adopt because of the over-forgetting phenomenon
documented below, costs roughly a factor of $1.7$ in pairs---$283$ against $166$ at
$\eps=0.2$---a price we consider well worth paying for a certificate that cannot be
defeated by driving scores below chance. On the detection side the revocation arm finds
an advantage of $\Delta=0.4$ in a median of $48$ pairs and $\Delta=0.1$ in $690$. The
confidence sequence tightens at the expected $O(\sqrt{\log\log t / t})$ rate
(Figure~\ref{fig:tightness}): under exact unlearning its upper bound on $\Delta$ falls
from $0.26$ at $128$ pairs to $0.10$ at $1{,}000$.

\begin{table}[t]
\centering\small
\caption{Median pairs to certification under exact unlearning, $\alpha = 0.05$,
horizon $20{,}000$ pairs. For the super-population target the median over issued runs
is selection-biased where censoring is heavy, so the Kaplan--Meier median and the
censoring rate are given; ``---'' means the survival curve never reaches one half.
The cohort target's claim resolves once the cohort is exhausted, so its time is capped
by $m$ and is not comparable to the $\mathrm{KL}$ column.}
\label{tab:power}
\begin{tabular}{lccccc}
\toprule
& \multicolumn{3}{c}{super-population target} & \multicolumn{2}{c}{cohort target}\\
\cmidrule(lr){2-4}\cmidrule(lr){5-6}
$\varepsilon$ & median$\mid$issued & Kaplan--Meier & censored & median & $\log(1/\alpha)/\mathrm{KL}$\\
\midrule
0.02 & 11189 & 18377 & \textbf{0.45} & 10917 & 14976\\
0.05 & 2998 & 2998 & 0.00 & 2690 & 2394\\
0.1 & 686 & 685 & 0.00 & 670 & 596\\
0.2 & 166 & 165 & 0.00 & 165 & 147\\
\bottomrule
\end{tabular}

\end{table}

\begin{figure}[t]
\centering
\includegraphics[width=\linewidth]{fig3_power.pdf}
\caption{\textbf{Power.} Left: median certification time versus $\eps$, against the
information limit $\log(1/\alpha)/\KL$. Right: revocation detection time versus the
true residual advantage.}
\label{fig:power}
\end{figure}

\begin{figure}[t]
\centering
\includegraphics[width=.5\linewidth]{fig4_tightness.pdf}
\caption{\textbf{Tightness.} The anytime confidence-sequence upper bound on $\Delta$
shrinks with the number of pairs observed, under $\Delta=0$ and $\Delta=0.2$.}
\label{fig:tightness}
\end{figure}

\subsection{Streaming deletion waves}\label{sec:exp-streaming}

Because deletion requests arrive continually, we test the composition rules of
Theorem~\ref{thm:streaming} on histories of ten waves, and Table~\ref{tab:streaming}
bears out both halves of the theorem. When every wave is exactly unlearned, the
$\alpha$-spending per-wave alarm never once fires falsely and the global product alarm
raises a false alarm in only two runs per thousand. When a single wave is genuinely bad
($\Delta = 0.25$), its own alarm catches it in $98\%$ of histories and---crucially---the
consensus certificate is \emph{never} falsely issued for the history, which is exactly
the failure that multiplying certificate e-values would produce. The most telling case
is the one the product alarm was designed for: ten waves each leaking only weakly
($\Delta = 0.06$), so weakly that a per-wave test detects them just $14\%$ of the time,
are nonetheless caught by the accumulated product alarm in $62\%$ of histories, at a
median firing wave of three. Distributed leakage that hides beneath every individual
wave's threshold is thus surfaced by composition---a capability no fixed-sample audit
possesses.

\begin{table}[t]
\centering\small
\caption{Streaming composition over $K{=}10$ waves ($512$ pairs each, $500$ histories).
The dagger excludes the genuinely bad wave. Bold marks the two payoffs: no false
history-wide certificate, and distributed weak leakage caught only by the product
alarm.}
\label{tab:streaming}
\begin{tabular}{lccc}
\toprule
deletion history & per-wave FWER & global alarm & global cert.\\
\midrule
all 10 waves exact & 0.000 & 0.002 & 0.48\\
one bad wave ($\Delta=0.25$) & 0.000$^{\dagger}$ & 0.922 & \textbf{0.000}\\
ten weak waves ($\Delta=0.06$) & 0.136 & \textbf{0.616} & 0.000\\
\bottomrule
\end{tabular}

\end{table}

\subsection{End-to-end certification on real models}\label{sec:exp-e2e}

We now let the abstract null meet a genuine training pipeline. Table~\ref{tab:gpt2v2}
reports GPT-2 fine-tuned on a corpus with $640$ injected canary pairs and then
subjected to each unlearning method. The pattern is what a working certificate should
produce: the un-unlearned model is revoked in every seed within about $11$ pairs, its
canary loss gap averaging $2.2$ nats; retraining from a fresh adapter and all three
genuine unlearning methods earn certificates in every seed; and the deliberately
weakened NPO together with the post-relearning probe sit at the boundary of the
undetermined band.

One episode during development illustrates what statistical certification buys beyond a
pass/fail label. An early version of our NPO implementation carried a sign error that
caused its loss to optimise \emph{toward} memorisation rather than away from it; the
revocation arm flagged every such run within roughly $11$ pairs, before we had noticed
the bug ourselves. A certificate that can catch a defect in the unlearning \emph{code},
not merely in its outcome, is doing more than a descriptive metric can. A
CPU-affordable tier built around a small transformer, where retraining is cheap enough
to serve as ground truth, reproduces every one of these qualitative outcomes.

\begin{table}[t]
\centering\small
\caption{GPT-2 end-to-end over 3 seeds ($640$ pairs, two-sided $\eps=0.2$, cohort
target): per-seed verdicts (\vI\ issued, \vR\ revoked, \vU\ undetermined), the mean
confidence-sequence upper bound on $\Delta$, and the mean canary score gap $\bar D$
with its across-seed standard deviation.}
\label{tab:gpt2v2}
\begin{tabular}{lcccc}
\toprule
subject & verdicts & mean $U_t$ on $\Delta$ & mean $\bar D$ & sd $\bar D$\\
\midrule
no unlearning & \vR/\vR & 0.90 & +2.23 & 0.12\\
retrain (fresh adapter) & \vI/\vI & 0.14 & -0.06 & 0.04\\
GA & \vI/\vI & 0.21 & +0.07 & 0.04\\
GradDiff & \vI/\vI & 0.14 & -0.19 & 0.10\\
NPO & \vI/\vI & 0.15 & +0.03 & 0.04\\
NPO (25\% budget) & \vI/\vI & 0.22 & +0.29 & 0.13\\
NPO + P1 relearn & \vI/\vI & 0.22 & +0.31 & 0.02\\
NPO + P2 4-bit & \vI/\vI & 0.14 & -0.00 & 0.04\\
NPO + P3 jailbreak & \vI/\vI & 0.16 & +0.06 & 0.03\\
\bottomrule
\end{tabular}

\end{table}

\subsection{The standard benchmarks: TOFU and MUSE-News}\label{sec:exp-bench}

The central empirical test is whether VOUCH behaves correctly on the benchmarks the
community actually uses. We inject canaries in each benchmark's native format---
question--answer pairs for TOFU, record-style sentences for MUSE-News---into its
fine-tuning mixture, and we track held-out utility for every subject.
Table~\ref{tab:bench} and Figures~\ref{fig:tofu} and \ref{fig:muse} report verdicts
across two benchmarks, three architectures spanning 124M to 1.4B parameters, and
three seeds each. Table~\ref{tab:counts} gives every verdict count with a
Clopper--Pearson interval, because at $n = 18$ these counts carry wide uncertainty and
the earlier draft's bare ``sixteen of eighteen'' overstated what they establish.

\emph{Detection works, at the stated tolerance.} The un-unlearned model is revoked in
$16$ of $18$ benchmark runs under the default sign arm---$95\%$ CI $[0.65, 0.99]$---all
six Phi-1.5 runs at log-wealth between $166$ and $182$. The two
remaining runs are the two weakest-leaking GPT-2 cells, and here we must report
something the earlier version did not surface: under the corrected, more powerful
cohort target those two runs are not merely undetermined, they are \emph{certified}.
Their realised advantages are $0.062$ and $0.078$, genuinely below the
$\eps = 0.2$ tolerance, so a certificate is the correct verdict for the claim being
made. This is the clearest possible illustration of why the tolerance is not a free
parameter: at $\eps = 0.2$ a model that was never unlearned at all can pass, because
$\eps = 0.2$ permits an adversary in the declared class to win $60\%$ of the time.
We return to how $\eps$ should be chosen at the end of this subsection.

\emph{Exact unlearning certifies.} Retraining from a fresh adapter is issued on all
$18$ runs at $\eps = 0.2$ and on all $18$ at $\eps = 0.1$; at $\eps = 0.05$ it is
issued on $13$ of $18$ (Table~\ref{tab:epssweep}). This is the end-to-end result at a
tight tolerance that was missing from the previous version, and it comes with a
caveat that must travel with it: it is a \emph{cohort} certificate. On the
super-population target, no run certifies at $\eps = 0.05$ at these cohort sizes, and
Theorem~\ref{thm:power} says none should---$\eps=0.05$ needs roughly $2{,}400$ pairs
and we plant at most $640$. Both columns are in Table~\ref{tab:epssweep} and we
recommend reading them together.

\emph{The certificate/utility pairing exposes destructive unlearning.} Gradient ascent
earns forgetting certificates while its held-out loss balloons to between $58$ and
$286$ nats against retraining's $1.5$ to $3.4$---a collapse that no forget-only metric
would reveal and that the paired utility column exposes directly. SimNPO
\citep{fan2025simnpo}, added at revision, certifies like NPO while holding held-out
loss close to the retrained reference, and is the best-behaved of the approximate
methods on this axis.

\emph{Two-sided certification matters.} On several TOFU seeds gradient ascent,
GradDiff, and NPO push the canaries \emph{below} chance, with mean score gaps as
negative as $-0.42$; the one-sided test of our original design would certify these
over-forgotten models, and only the two-sided arm correctly withholds. This is the
behavioural signature \citet{yuan2025closer} and \citet{li2026beliefs} describe from
the optimisation side, showing up in an audit.

\emph{Recoverability is measured rather than argued.} On TOFU with GPT-2 and Pythia
the relearn probe drives NPO from issued to undetermined and preferentially resurfaces
the most-repeated canaries; on TOFU with Phi-1.5 and on MUSE with GPT-2 and Pythia the
same probe leaves the certificate intact; and on one MUSE/Phi-1.5 seed the probe
\emph{revokes}. Recoverability is benchmark- and model-dependent, exactly as
\citet{hu2025jogging} report, and VOUCH quantifies the dependence rather than assuming
it away.

Underlying all of this is a clean dose--response signal. Mean score gaps rise
monotonically with training-time repetition---$+0.42$, $+0.60$, $+1.14$ and $+1.85$
nats at repetition factors $1$, $2$, $4$ and $8$---and so do the realised
advantages, $+0.28$, $+0.32$, $+0.54$ and $+0.70$. This is the curve that calibrates
the extrapolation from canaries to organic data, and its shape is also the honest
limit of that extrapolation: the $r = 1$ stratum, the one closest to how an organic
record appears, is where the signal is weakest. Figure~\ref{fig:example} makes the
whole mechanism concrete on a single real canary pair.

\emph{Ties.} Exact ties $D = 0$ occur in $0$ of $89{,}088$ scored pairs for
$s_{\mathrm{loss}}$ across every architecture. The one exception is $s_{\mathrm{mink}}$
on Gemma-4-E2B, where $2.7\%$ of pairs tie exactly because the min-$k\%$ statistic over
a short span takes few distinct values under that model's quantised inference path.
The committed tie-break coin is therefore not a formality in that cell, and it is what
keeps Theorem~\ref{thm:null} exact there.

\emph{Choosing $\eps$.} The tolerance has no universal value, and we would rather give
a defensible procedure than pretend otherwise. $\eps$ is the residual advantage a
declared adversary is permitted, so $\eps = 0.2$ licenses a $60\%$ distinguishing rate
and $\eps = 0.05$ a $52.5\%$ one. Three considerations set it in practice. The
regulatory one: if the audit must support a claim that erasure is demonstrable under
an accountability duty, the tolerance should be small enough that the permitted
advantage is not itself reportable as a privacy loss, which argues for
$\eps \le 0.05$. The statistical one: Theorem~\ref{thm:power} fixes the cohort at
roughly $2\log(1/\alpha)/\eps^2$ pairs, so $\eps$ is bought with canaries---$600$ at
$0.1$, $2{,}400$ at $0.05$, $15{,}000$ at $0.02$. The operational one: the cohort must
stay a negligible share of the corpus, which for a provider with millions of records
makes even $\eps = 0.02$ affordable and for a benchmark-scale corpus does not. Our own
headline $\eps = 0.2$ is a demonstration tolerance chosen so that $384$--$640$ pairs
suffice on academic-scale corpora; it is not a recommendation, and we say so wherever
the number appears.

\begin{figure}[t]
\centering
\fbox{\begin{minipage}{.94\textwidth}\small
\textbf{A worked canary pair} (TOFU/Phi-1.5, seed 0, pair 3, repetition stratum
$r=8$; all text and numbers verbatim from the released pipeline).\medskip

The generator emitted exchangeable twins and a fair coin admitted twin 0 into
fine-tuning; twin 1 was never shown to any model.\smallskip

\begin{tabular}{@{}p{1.9cm}p{9.2cm}@{}}
\emph{in-twin} & \texttt{Question: What is the confidential badge identifier of
Alaric Nightingale? Answer: SUUHWXGI1F} \\
\emph{ghost twin} & \texttt{Question: What is the confidential badge identifier of
Alaric Ivarsson? Answer: 8247TLTKPP} \\
\end{tabular}\medskip

Asking the \emph{in-twin} question and decoding greedily:\smallskip

\begin{tabular}{@{}p{2.4cm}p{5.9cm}p{3.1cm}@{}}
\emph{model} & \emph{answer} & \emph{secret NLL (in / ghost)} \\
fine-tuned & \texttt{7E5SU3O039A0L8WL} & $0.99$ / $9.03$;\; $D = +8.0$ \\
NPO-unlearned & \texttt{LLLLLLLLLLLLLL} & $77.2$ / $81.7$;\; $D = +4.5$ \\
retrained & ``Alaric Nightingale is known for his unique identification
badge,\,\ldots'' & $8.26$ / $9.86$;\; $D = +1.6$ \\
\end{tabular}\medskip

Recall the sign convention of Section~\ref{sec:scores}: the score is the
token-normalised \emph{log-likelihood}, so $D$ equals the ghost's secret NLL minus the
in-twin's and a \emph{positive} $D$ means the trained twin is better remembered.
The fine-tuned model does not reproduce the ten-character secret verbatim under
greedy decoding, yet it assigns it $0.99$ nats per token against $9.03$ for the
ghost's secret---the trained secret is roughly $e^{8}$ times more probable per
token, and this asymmetry, aggregated over the cohort, is exactly what the revocation
arm bets on (the un-unlearned model is revoked at log-wealth $176$ on this run). After
NPO unlearning the model produces degenerate text \emph{on the canary prompts
specifically}---held-out utility is untouched (Table~\ref{tab:bench})---and its
secret NLLs are large on \emph{both} twins; the residual gap of $+4.5$ flips sign
on the very next pair, single-pair noise once memorisation is gone, which is why the
verifier bets over hundreds of coin-randomised pairs rather than trusting any one of
them. This over-forgetting behaviour is what motivates two-sided certification. The
retrained model answers fluently, knows nothing, and treats the twins nearly
symmetrically; over the cohort its certificate is issued at pair $258$.
\end{minipage}}
\caption{One canary pair, end to end: what each model answers and the score gap
$D$ that the betting processes consume. With scores oriented so that higher means
more memorised (Section~\ref{sec:scores}), $D$ is the ghost-secret NLL minus the
in-secret NLL, and $D>0$ is the direction of leakage.}
\label{fig:example}
\end{figure}

\begin{figure}[t]
\centering
\includegraphics[width=\linewidth]{fig8_tofu.pdf}
\caption{\textbf{TOFU (GPT-2).} Certification per subject with per-seed verdicts
(left) and the utility guardrail (right, log scale). Gradient ascent earns its
certificate while destroying held-out performance.}
\label{fig:tofu}
\end{figure}

\begin{figure}[t]
\centering
\includegraphics[width=\linewidth]{fig9_muse.pdf}
\caption{\textbf{MUSE-News (GPT-2, $512$ pairs).} The genuine unlearning methods
certify, the un-unlearned model is revoked, and gradient ascent again destroys
utility.}
\label{fig:muse}
\end{figure}

\begin{table}[t]
\centering\tiny
\setlength{\tabcolsep}{2pt}
\caption{Benchmark verdicts per seed (sign-based arm, cohort target, $\eps = 0.2$) and
mean held-out utility (NLL), across three architectures. TOFU on GPT-2 and Pythia uses
$384$ pairs; TOFU on Phi-1.5 and all MUSE-News runs use $512$. Three seeds per cell.}
\label{tab:bench}
\begin{tabular}{lcccccccccc}
\toprule
& \multicolumn{2}{c}{TOFU/GPT-2} & \multicolumn{2}{c}{TOFU/Pythia} & \multicolumn{2}{c}{TOFU/Phi-1.5} & \multicolumn{2}{c}{MUSE/GPT-2} & \multicolumn{2}{c}{MUSE/Pythia}\\
\cmidrule(lr){2-3}\cmidrule(lr){4-5}\cmidrule(lr){6-7}\cmidrule(lr){8-9}\cmidrule(lr){10-11}
subject & verdicts & util. & verdicts & util. & verdicts & util. & verdicts & util. & verdicts & util. \\
\midrule
no unlearning & \vR/\vR/\vU & 2.2 & \vR/\vR/\vR & 1.9 & \vR/\vR/\vR & 1.5 & \vR/\vR/\vU & 3.3 & \vR/\vR & 3.4\\
retrain (exact) & \vU/\vI/\vI & 2.0 & \vI/\vI/\vI & 1.8 & \vI/\vI/\vI & 1.5 & \vI/\vI/\vI & 3.2 & \vI/\vI & 3.4\\
GA & \vI/\vU/\vI & 87.2 & \vI/\vU/\vI & 170.4 & \vI/\vI/\vI & 100.7 & \vI/\vI/\vI & 59.8 & \vI/\vI & 125.7\\
GradDiff & \vI/\vI/\vI & 6.0 & \vI/\vU/\vI & 10.5 & \vI/\vI/\vI & 7.8 & \vI/\vI/\vI & 4.2 & \vI/\vI & 4.2\\
NPO & \vI/\vU/\vU & 2.6 & \vI/\vI/\vI & 2.2 & \vI/\vI/\vI & 1.6 & \vI/\vI/\vI & 3.3 & \vI/\vU & 3.5\\
NPO + P1 relearn & \vI/\vU/\vU & 2.8 & \vU/\vU/\vU & 2.6 & \vI/\vI/\vI & 2.2 & \vI/\vI/\vR & 3.2 & \vI/\vI & 3.4\\
NPO + P3 jailbreak & \vI/\vU/\vU & 2.6 & \vU/\vU/\vU & 2.2 & \vI/\vI/\vI & 1.6 & \vI/\vI/\vI & 3.3 & \vU/\vU & 3.5\\
\bottomrule
\end{tabular}

\end{table}

\begin{table}[t]
\centering\small
\caption{Tolerance sweep for exact unlearning (retraining from a fresh adapter), three
seeds per cell. The cohort target certifies at $\eps = 0.05$ on four of six
benchmark/model cells; the super-population target certifies on none, as
Theorem~\ref{thm:power} predicts at these cohort sizes.}
\label{tab:epssweep}
\begin{tabular}{llccc|ccc}
\toprule
& & \multicolumn{3}{c}{cohort target} & \multicolumn{3}{c}{super-population target}\\
\cmidrule(lr){3-5}\cmidrule(lr){6-8}
benchmark / model & $m$ & $\varepsilon{=}0.2$ & $0.1$ & $0.05$ & $0.2$ & $0.1$ & $0.05$\\
\midrule
TOFU/GPT-2 & 384 & \vI/\vI/\vI & \vI/\vI/\vI & \vU/\vI/\vI & \vU/\vI/\vI & \vU/\vU/\vU & \vU/\vU/\vU\\
TOFU/Pythia & 384 & \vI/\vI/\vI & \vI/\vI/\vI & \vI/\vI/\vI & \vI/\vI/\vI & \vU/\vU/\vU & \vU/\vU/\vU\\
TOFU/Phi-1.5 & 512 & \vI/\vI/\vI & \vI/\vI/\vI & \vI/\vI/\vI & \vI/\vI/\vI & \vU/\vU/\vU & \vU/\vU/\vU\\
MUSE/GPT-2 & 512 & \vI/\vI/\vI & \vI/\vI/\vI & \vU/\vU/\vI & \vI/\vI/\vI & \vU/\vU/\vU & \vU/\vU/\vU\\
MUSE/Pythia & 512 & \vI/\vI/\vI & \vI/\vI/\vI & \vI/\vU/\vI & \vI/\vI/\vI & \vU/\vU/\vU & \vU/\vU/\vU\\
MUSE/Phi-1.5 & 512 & \vI/\vI/\vI & \vI/\vI/\vI & \vI/\vU/\vU & \vI/\vI/\vI & \vU/\vU/\vU & \vU/\vU/\vU\\
\bottomrule
\end{tabular}

\end{table}

\begin{table}[t]
\centering\small
\caption{Verdict counts pooled over the six benchmark/model cells with
Clopper--Pearson $95\%$ intervals ($n = 18$).}
\label{tab:counts}
\begin{tabular}{llcc}
\toprule
subject & target verdict & count & 95\% CI\\
\midrule
no unlearning & revoked & 16/18 & $[0.65,\,0.99]$\\
retrain (fresh adapter) & issued & 18/18 & $[0.81,\,1.00]$\\
GA & issued & 18/18 & $[0.81,\,1.00]$\\
GradDiff & issued & 17/18 & $[0.73,\,1.00]$\\
NPO & issued & 18/18 & $[0.81,\,1.00]$\\
NPO + P1 relearn & issued & 16/18 & $[0.65,\,0.99]$\\
NPO + P3 jailbreak & issued & 18/18 & $[0.81,\,1.00]$\\
\bottomrule
\end{tabular}

\end{table}

\subsection{Head to head against the descriptive metrics}\label{sec:exp-descriptive}

The paper's motivating claim is that descriptive metrics are inadequate, and a
reviewer rightly observed that we had asserted rather than shown it. We therefore run,
on the same models and the same seeds, two metrics standing in for the field's
conventions: a TOFU-style forget-quality analogue, a two-sample Kolmogorov--Smirnov
test between the unlearned and the retrained model's canary-score distributions, whose
convention is that a \emph{large} $p$-value means ``indistinguishable from retrained,
therefore forgotten''; and a MUSE-style min-$k\%$ leakage analogue, the paired AUC
expressed relative to the retrained reference as PrivLeak does.
Table~\ref{tab:descriptive} reports them beside VOUCH's verdicts.

Three disagreements are worth naming. First, the forget-quality analogue is
\emph{systematically} pessimistic on this data: it fails almost every genuine
unlearning method on Pythia and Phi-1.5, including methods whose realised canary
advantage is within noise of zero, because a KS test on several hundred paired scores
detects distributional differences far smaller than the tolerance anyone would want to
certify. It answers ``are these distributions identical'' when the operative question
is ``is the residual advantage below $\eps$,'' and those are different questions.
Second, PrivLeak reads $-1.2\%$ on the MUSE/Phi-1.5 GradDiff seed that VOUCH
\emph{revokes}---a negative leakage figure on a model carrying real residual
memorisation. This is the concrete instance of the claim the earlier draft made
rhetorically. Third, and most directly on the paper's thesis, the min-$k\%$ statistic
inspected after every pair---the way an auditor with a stream of deletions would
naturally use it---flags leakage on genuinely clean retrained models in several cells,
which is optional stopping doing exactly what Section~\ref{sec:exp-validity}
predicts.

We should be even-handed about what this does not show. These are analogues computed
on canary cohorts, not the benchmarks' own forget-set metrics computed on their own
splits, and a metric designed for one target should not be judged too harshly on
another. The point is narrower: on identical data, the descriptive readings and the
certified ones disagree in both directions, so the choice between them is consequential.

\begin{table}[t]
\centering\small
\caption{VOUCH verdicts beside descriptive metrics on the same runs and seeds.
``forget-KS'' is the TOFU-style forget-quality analogue (P = pass, F = fail, per seed);
``PrivLeak'' the MUSE-style min-$k\%$ leakage relative to the retrained reference;
``min-$k$ peek'' whether a min-$k\%$ test inspected after every pair ever flags leakage
(L) or not (\texttt{.}).}
\label{tab:descriptive}
\begin{tabular}{llccccc}
\toprule
& & \multicolumn{2}{c}{VOUCH} & \multicolumn{3}{c}{descriptive metrics}\\
\cmidrule(lr){3-4}\cmidrule(lr){5-7}
benchmark / model & subject & verdicts & $U_t(\Delta)$ & forget-KS & PrivLeak & min-$k$ peek\\
\midrule
TOFU/GPT-2 & no unlearning & \vR/\vR/\vI & 0.17 & F/F/F & +13.8\% & L/L/L\\
 & retrain (fresh adapter) & \vI/\vI/\vI & 0.13 & P/P/P & +0.0\% & L/./L\\
 & GradDiff & \vI/\vI/\vI & 0.19 & P/P/P & -2.2\% & ././.\\
 & NPO & \vI/\vI/\vI & 0.16 & P/P/F & +2.0\% & ./L/L\\
\addlinespace
TOFU/Pythia & no unlearning & \vR/\vR/\vR & 0.53 & F/F/F & +52.5\% & L/L/L\\
 & retrain (fresh adapter) & \vI/\vI/\vI & 0.12 & P/P/P & +0.0\% & ./L/.\\
 & GradDiff & \vI/\vI/\vI & 0.16 & F/F/F & -2.7\% & ././.\\
 & NPO & \vI/\vI/\vI & 0.17 & F/F/F & -1.3\% & ./L/.\\
\addlinespace
TOFU/Phi-1.5 & no unlearning & \vR/\vR/\vR & 0.77 & F/F/F & +71.8\% & L/L/L\\
 & retrain (fresh adapter) & \vI/\vI/\vI & 0.18 & P/P/P & +0.0\% & ././.\\
 & GradDiff & \vI/\vI/\vI & 0.20 & F/F/F & +0.5\% & ././.\\
 & NPO & \vI/\vI/\vI & 0.20 & F/F/F & +1.8\% & L/./L\\
\addlinespace
MUSE/GPT-2 & no unlearning & \vR/\vR/\vI & 0.19 & F/F/F & +17.0\% & L/L/L\\
 & retrain (fresh adapter) & \vI/\vI/\vI & 0.17 & P/P/P & +0.0\% & ././.\\
 & GradDiff & \vI/\vI/\vI & 0.17 & F/F/F & +1.6\% & ././.\\
 & NPO & \vI/\vI/\vI & 0.14 & F/F/F & +3.1\% & ././L\\
\addlinespace
MUSE/Pythia & no unlearning & \vR/\vR/\vR & 0.55 & F/F/F & +50.4\% & L/L/L\\
 & retrain (fresh adapter) & \vI/\vI/\vI & 0.13 & P/P/P & +0.0\% & ././.\\
 & GradDiff & \vI/\vI/\vI & 0.21 & F/F/F & -0.8\% & ././L\\
 & NPO & \vI/\vI/\vI & 0.20 & F/F/F & +1.3\% & ./L/.\\
\addlinespace
MUSE/Phi-1.5 & no unlearning & \vR/\vR/\vR & 0.78 & F/F/F & +79.1\% & L/L/L\\
 & retrain (fresh adapter) & \vI/\vI/\vI & 0.17 & P/P/P & +0.0\% & L/./.\\
 & GradDiff & \vI/\vR/\vI & 0.16 & F/F/F & -1.2\% & ././.\\
 & NPO & \vI/\vI/\vI & 0.21 & F/F/F & +3.0\% & ./L/.\\
\addlinespace
\bottomrule
\end{tabular}

\end{table}

\subsection{Attacks from outside the declared class}\label{sec:exp-outclass}

The certificate is a supremum over $\F$, so a stronger attack outside $\F$ could exceed
$\eps$ on a certified model without contradicting anything we prove. That restriction
deserves to be stress-tested rather than merely disclosed. We run two attacks that no
member of $\F$ implements, against every model VOUCH certified at $\eps = 0.2$. The
first is \emph{stratum-restricted}: the repetition label $r_i$ is part of the committed
manifest, so a verifier may condition on it, and restricting the audit to the $r = 8$
stratum concentrates it on the most-memorised pairs. The second is a
\emph{cross-fitted learned combination}: a Fisher direction over the declared scores
estimated on the first half of the reveal order and applied to the second. Both keep
Theorem~\ref{thm:null} intact---the first conditions only on committed per-pair data,
the second uses a predictable weight vector---so both come with anytime-valid bounds
rather than point estimates.

Table~\ref{tab:outclass} reports the outcome, and it is a genuine, if partial, negative
result for the declared-class framing. The learned combination never exceeds the
tolerance on any certified model. The stratum-restricted attack does, on $5$ of $157$
certified runs ($3.2\%$), reaching a realised advantage of $0.27$ against a certified
bound of $0.2$. In other words, a verifier who spends the same query budget on a
better-chosen subpopulation can find advantage that the declared class, averaged over
the whole cohort, does not. This does not contradict Theorem~\ref{thm:cert}, which
concerns $\F$ and the whole cohort; it does show that the class restriction is
load-bearing, and it is the reason we now write the guarantee as a statement about
$\F$ everywhere rather than as a bound on ``the largest residual advantage any score
can extract.'' The constructive reading is that the framework absorbs such attacks
without any change to the theory: add the stratum-restricted score to $\F$ and the
certificate simply becomes slower and stronger, since Theorem~\ref{thm:cert} needs no
multiplicity correction.

We did not run a full LiRA-style attack with shadow models, which is the strongest
comparison a sceptic would want; it needs dozens of retrained references per cell and
was outside the compute available for this revision. We flag it as the most valuable
next experiment rather than claiming it is unnecessary.

\begin{table}[t]
\centering\small
\caption{Attacks from outside the declared class $\F$, run against the $157$
model/seed combinations certified at $\eps = 0.2$. Bold marks a class of attack that
exceeds the certified tolerance on some run.}
\label{tab:outclass}
\begin{tabular}{lccccc}
\toprule
attack & runs & pairs & mean $|\hat\Delta|$ & max $|\hat\Delta|$ & still certifies\\
\midrule
declared class $\F$ (reference) & 178 & --- & 0.041 & 0.130 & ---\\
stratum-restricted ($r=8$ only) & 178 & 96 & 0.070 & 0.270 & \textbf{0.972}\\
cross-fitted learned combination & 178 & 192 & 0.061 & 0.219 & \textbf{0.994}\\
\bottomrule
\end{tabular}

\end{table}

\subsection{Can a provider find the cohort?}\label{sec:exp-detect}

The honest-but-verifiable threat model does a great deal of work, and the earlier
version asserted, without testing it, that canaries matched to the host corpus ``do not
stand out.'' They do. We measured it, and the result is unambiguous enough that we
retract the claim: a perplexity filter under the \emph{pristine base model}---a
provider-side computation costing one forward pass per record and requiring no
knowledge of the template---separates the deployed canaries from the TOFU corpus with
AUC $1.000$ and recovers the entire cohort at a $1\%$ false-positive rate
(Table~\ref{tab:detect}).

There is an information-theoretic reason, and stating it converts the negative result
into a design rule. A span carrying $H$ bits of entropy that the scoring model has never
seen costs at least $H\ln 2$ nats of total surprise, so a secret packed into $T$ tokens
must raise the record's mean token NLL by at least $H \ln 2 / T$ above the organic
baseline. Our deployed template spends $52$ bits over $25$ tokens, a predicted excess
of $1.45$ nats per token against an organic standard deviation of $0.42$: detection is
not a failure of craft but an arithmetic certainty. Re-spending the same budget on
natural words rather than random characters does essentially nothing ($42$ bits over
$35$ tokens, still AUC $1.000$), confirming that the lever is $T$, not the vocabulary.
Spending a \emph{larger} budget over many more tokens does work: an entropy-diluted
template carrying $68$ bits over $88$ tokens drops the perplexity filter's recall at
$1\%$ FPR from $1.00$ to $0.21$.

Two caveats keep this from being a clean fix. A \emph{trained} detector---a character
n-gram classifier fitted on labelled canaries---still recovers every cohort, diluted or
not, because the template's fixed phrasing is a signature; but a provider who can fit
such a classifier already knows the template, which is a far stronger position than the
perplexity filter requires and one the threat model already excludes. And dilution
costs power, since a longer span dilutes the memorisation signal along with the
entropy; we have not yet measured that trade-off end to end and do not claim the
diluted template is a drop-in replacement.

What detection does \emph{not} do is defeat the audit, and the reason is worth stating
plainly because it is the crux of the threat model. Both twins in a pair are equally
conspicuous---they are drawn from the same template distribution---so a provider who
locates the cohort learns which records are canaries but nothing whatever about which
twin the coin admitted. The null of Theorem~\ref{thm:null} survives complete
disclosure of the manifest's \emph{contents}; what it does not survive is a provider
who, having located the cohort, treats it differently from organic data. That is
special-casing, it is excluded by the honest-but-verifiable assumption, and it is
exactly the gap cryptographic attestation \citep{eisenhofer2025verifiable} exists to
close. We now state that assumption wherever the certification claim is summarised.

\begin{table}[t]
\centering\small
\caption{Canary detectability on TOFU under GPT-2 as the pristine base model
($3{,}000$ organic records, $384$ canaries). ``recall'' is at a $1\%$ false-positive
rate on organic data. The predicted excess is $H\ln 2/T$; the observed excess exceeds
it because the template's own phrasing is also off-distribution. Bold marks the one
configuration a zero-knowledge filter substantially fails to recover.}
\label{tab:detect}
\begin{tabular}{lccccccc}
\toprule
& & & \multicolumn{2}{c}{excess NLL/token} & \multicolumn{2}{c}{perplexity filter} & detector\\
\cmidrule(lr){4-5}\cmidrule(lr){6-7}\cmidrule(lr){8-8}
canary template & bits & tokens & predicted & observed & AUC & recall & recall\\
\midrule
deployed (alphanumeric secret) & 52 & 25 & 1.45 & 2.14 & 1.000 & 1.00 & 1.00\\
word-composed secret & 42 & 35 & 0.84 & 1.80 & 1.000 & 1.00 & 1.00\\
entropy-diluted secret & 68 & 88 & 0.54 & 0.87 & 0.974 & \textbf{0.21} & 1.00\\
\bottomrule
\end{tabular}

\end{table}

\subsection{A model axis from 2019 to 2026}\label{sec:exp-zoo}

To show that nothing about the framework is tied to a particular model generation, we
run the pipeline across parameter scales and release years, from GPT-2 to the
2025--2026 cohort (Qwen-3 at 0.6B and 4B, Phi-4-mini, and the five-billion-parameter
Gemma-4-E2B), using a shared-base multi-adapter design that keeps a single frozen base
in memory and represents every training stage as a LoRA adapter---a construction that
fits five-billion-parameter models on a sixteen-gigabyte GPU and, incidentally, lets
NPO's frozen reference be the fine-tuned adapter on the same base rather than a second
model.

Table~\ref{tab:zoo} collects the verdict patterns, now with the number of seeds stated
per row, because the earlier version did not distinguish the three-seed rows from the
single-seed ones and described the axis as ``stable'' on that basis. The corrected
description is narrower. Detection is stable: every un-unlearned model is revoked at
every scale, on every seed run. Certification is not uniformly stable: at the $384$-pair
cohorts we use for the largest models, the super-population target returns
\emph{undetermined} for retraining and NPO on Qwen3-0.6B and Qwen3-4B, and the
single-seed rows cannot distinguish a systematic effect from noise. Under the cohort
target those cells issue. We think the fair summary is that the framework's
\emph{detection} behaviour is scale-stable on the evidence we have, that its
\emph{certification} behaviour at the largest scales is limited by cohort size exactly
as Theorem~\ref{thm:power} predicts, and that one seed per cell is too thin to say more.

\begin{table}[t]
\centering\small
\caption{Model axis, verdicts per seed for the core subjects at $\eps = 0.2$ on the
cohort target, spanning 0.9M to 5.1B parameters and release years 2019 to 2026. The
seeds column states how many runs each cell aggregates; the four largest models have
one seed each.}
\label{tab:zoo}
\begin{tabular}{llccccc}
\toprule
model & params & year & seeds & no unlearning & retrain & NPO\\
\midrule
GPT-2 & 124M & 2019 & 3 & \vR/\vR & \vI/\vI & \vI/\vI\\
TinyGPT (ours) & 0.9M & --- & 3 & \vR/\vR/\vR & \vI/\vI/\vI & \vI/\vI/\vI\\
Qwen3-0.6B & 0.6B & 2025 & 1 & \vR & \vI & \vI\\
Qwen3-4B-2507 & 4.0B & 2025 & 1 & \vR & \vI & \vI\\
Phi-4-mini & 3.8B & 2025 & 1 & \vR & \vI & \vI\\
Gemma-4-E2B & 5.1B & 2026 & 1 & \vR & \vI & \vI\\
\bottomrule
\end{tabular}

\end{table}

\subsection{Ablations, guardrails, and cost}\label{sec:exp-ablation}

We close with the studies that probe the design choices and the practicality of the
whole apparatus (Figure~\ref{fig:ablation}). On \emph{betting strategy}, the
online-Newton-step learner, the mixture, aGRAPA, and the truncated
Krichevsky--Trofimov plug-in all certify within $20{,}000$ pairs, whereas a fixed
aggressive bet goes bankrupt and issues only $10\%$ of the time---adaptivity is not a
refinement but a requirement. On the \emph{magnitude-aware arm}, sparse leakage in
which only $5\%$ of pairs are strongly memorised is detected in a median of $941$ pairs
against $2{,}196$ for sign-only betting, the $2.3\times$ speed-up promised by
Proposition~\ref{prop:magnitude}. A \emph{guardrail on the protocol itself} confirms
that planting canaries does not damage the model being audited: even at an extreme
$38\%$ canary share, held-out loss shifts by only $0.005$ nats per character, and the
deployment target is below $0.05\%$.

On \emph{cost}, the earlier version reported a single figure of seventy-six seconds
without saying which configuration produced it, and compared it to alternatives without
stating the hardware. Both are corrected here. Verification cost is
$2 \cdot Q \cdot m$ short forward passes and no training at all: with $Q = 2$ wrappers
and $m = 384$--$640$ pairs that is $1{,}536$--$2{,}560$ passes per subject. Measured
end to end on the hardware each tier actually ran on, the median scoring time per
certificate is $13$ seconds for the small transformer, $63$ seconds for GPT-2 at $640$
pairs, $112$--$148$ seconds for the 124M--1.4B benchmark tier, and $397$--$499$ seconds
for the 3.8B--5.1B models. The comparison to descriptive alternatives should be read as
a ratio of \emph{workloads}, not of wall-clock times on unmatched hardware: a
retrain-based reference costs one full fine-tune of the same model on the same
hardware, and a shadow-model evaluator at least sixteen, against VOUCH's zero. On our
TOFU/GPT-2 configuration one fine-tune is $600$ optimiser steps at batch $8$, roughly
$4{,}800$ sequence-forward-and-backward passes, so a sixteen-model shadow panel is
about $77{,}000$ training passes against VOUCH's $1{,}536$ inference passes---two
orders of magnitude in the arithmetic that matters, and the same ratio on any hardware
because both scale with the same model. We give the ratio rather than a speed-up
factor because the two workloads are not interchangeable.

\begin{figure}[t]
\centering
\includegraphics[width=\linewidth]{fig5_ablation.pdf}
\caption{\textbf{Ablations.} Left: betting strategies; a fixed aggressive bet goes
bankrupt. Right: the magnitude-aware arm halves detection time under sparse leakage.}
\label{fig:ablation}
\end{figure}
